\section{第二条主线：模型能力和系统工程不是替代关系}

\subsection{从 LM 到 LM system，再到 Agent system}
讨论中很有价值的一点，是嘉宾们没有把 ``模型派'' 和 ``系统派'' 对立起来。相反，他们把今天的 Agent 系统看成 LM system 的自然延伸：模型之外还需要 infra、memory、定制化 workflow、交互界面与工具封装。换句话说，强模型并不会让系统工程消失，只会改变系统工程的重点。

\textit{``它不只是一个模型在做服务，而是以大模型为核心的一套系统在服务给我们。''}

\begin{importantbox}{为什么 Cursor / Claude Code 会成为高频例子}
这些产品的吸引力并不只来自模型本身的 coding 能力，还来自整套交互系统：
\begin{itemize}
  \item 编辑器上下文如何注入；
  \item 文件与命令权限如何管理；
  \item 工具调用结果如何反馈给模型；
  \item 用户如何在关键节点接管流程。
\end{itemize}
这说明 ``好的 Agent 产品'' 既是模型问题，也是产品系统设计问题。
\end{importantbox}

\subsection{Memory、workflow、tool space 的边界在移动}
面向 2024 年的 prompting agent，很多能力主要依靠外部 workflow 实现；到了 2025 年，嘉宾们观察到越来越多能力开始往模型内部迁移，包括 coding 习惯、computer use 策略、某些 reasoning 模式，甚至部分 memory 与 tool use 能力。与此同时，外部系统并没有消失，而是在处理新的问题：更复杂的权限控制、更细的用户个性化、更长时间的状态管理。

\begin{table}[H]
\centering
\begin{tabular}{p{3.0cm}p{4.3cm}p{4.3cm}}
\toprule
\textbf{系统层} & \textbf{2024 年典型做法} & \textbf{2025 年讨论中的变化} \\
\midrule
模型内部 & 通用聊天、少量 tool use & 开始内化 coding、GUI、multi-agent、reasoning \\
Memory & 外部检索或临时缓存 & 讨论如何部分内化，但长期状态仍依赖系统 \\
Workflow & Prompt 编排为主 & 逐步转向能力更强的自主探索与策略学习 \\
Tool space & 少量固定工具 & 行动空间变宽，环境更复杂，约束更关键 \\
产品界面 & 简单 chat UI & 编辑器、移动端、企业系统入口深度耦合 \\
\bottomrule
\end{tabular}
\caption{Panel 中隐含的一条判断：模型与系统的边界在不断重划}
\end{table}

\subsection{能力增强之后，系统会承接 ``最后一公里''}
嘉宾们对下一阶段的判断很一致：当模型内部的 agentic 能力增强到一定程度后，研究重点会再次回到系统外层。原因并不复杂，强模型只能提升上限，但企业应用需要的是稳定性、权限隔离、个性化、安全与部署接口。这些问题即使在模型很强时，也必须由系统来承接。

\begin{warningbox}{把所有能力都内化进模型并不现实}
如果 environment、memory、权限、外部 API 全都靠模型 ``脑补''，系统就会失去可控性。嘉宾们并没有主张彻底内化，而是强调 ``哪里是短板，就优先补哪里''。这是一种非常工程化的判断，而不是教条式的 architecture 选择。
\end{warningbox}

\subsection{本章小结}
Agent 的发展不是 ``模型取代系统''，而是两者交替成为短板。2025 年的研究重点更偏向把能力注入模型内部；但一旦能力增强，memory、tooling、workflow、权限和产品化接口会重新成为瓶颈。这个来回摆动，本身就是 Agent 系统成熟的轨迹。

